% =============================================================================
% BM PREDICT-FCT: Forecast Methodology
% =============================================================================
% Category: PREDICT
% Version: 1.0
% Created: 2026-01-30
% Dependencies: V (CORE-WHEN), BBB (CORE-WHERE), AN (METHOD-LLMMC)
% Primary Chapter: N/A (Cross-cutting methodology)
% =============================================================================

\chapter{Forecast Methodology}
\label{app:forecast-methodology}

% -----------------------------------------------------------------------------
% METADATA BLOCK
% -----------------------------------------------------------------------------
\begin{tcolorbox}[colback=blue!5!white,colframe=blue!75!black,title=Appendix Metadata]
\begin{tabular}{ll}
\textbf{Code:} & BM \\
\textbf{Category:} & PREDICT-FCT \\
\textbf{Title:} & Forecast Methodology \\
\textbf{Version:} & 1.0 \\
\textbf{Created:} & 2026-01-30 \\
\textbf{Status:} & Active \\
\end{tabular}
\end{tcolorbox}

% -----------------------------------------------------------------------------
% CROSS-REFERENCE MAP
% -----------------------------------------------------------------------------
\begin{tcolorbox}[colback=gray!5!white,colframe=gray!75!black,title=Cross-Reference Map]
\textbf{Dependencies:}
\begin{itemize}
    \item Appendix V (CORE-WHEN): Context dimensions $\Psi$
    \item Appendix BBB (CORE-WHERE): Parameter sources
    \item Appendix AN (METHOD-LLMMC): LLM Monte Carlo estimation
    \item Appendix MS (REF-MODELSPACE): Theory catalog
\end{itemize}

\textbf{Related Appendices:}
\begin{itemize}
    \item Appendix S (PREDICT-FALSIFY): Falsifiable predictions
    \item Appendix EEE (METHOD-DESIGN): Model design workflow
    \item Appendix GGG (METHOD-CONFIG): Configuration defaults
\end{itemize}
\end{tcolorbox}

% -----------------------------------------------------------------------------
% ABSTRACT
% -----------------------------------------------------------------------------
\begin{abstract}
This appendix documents the EBF forecast methodology for systematic prediction tracking, probability updating, and calibration measurement. It covers political referenda, geopolitical events, economic forecasts, and behavioral predictions. The methodology emphasizes: (1) precise resolution criteria defined \textit{ex ante}, (2) probability history tracking for Bayesian updating, (3) conditional scenario analysis, and (4) post-hoc calibration via Brier scores and log loss. The framework integrates with the EBF 10C dimensions and context analysis ($\Psi$-framework) to ensure forecasts are grounded in behavioral economics principles.
\end{abstract}

% -----------------------------------------------------------------------------
% QUICK REFERENCE
% -----------------------------------------------------------------------------
\begin{tcolorbox}[colback=green!5!white,colframe=green!75!black,title=Quick Reference]
\textbf{Key Concepts in this Appendix:}
\begin{itemize}
    \item \textbf{Forecast Registry:} Primary database \texttt{forecast-registry.yaml}
    \item \textbf{Four-Database Architecture:} Forecast $\rightarrow$ Session $\rightarrow$ Model $\rightarrow$ Output
    \item \textbf{Brier Score:} $BS = (P - O)^2$ where $O \in \{0,1\}$
    \item \textbf{Probability History:} Time-series $\mathcal{H} = \{(t_i, P_i, \text{CI}_i, R_i)\}$
    \item \textbf{Resolution Criteria:} Ex-ante definition of event occurrence
    \item \textbf{Conditional Scenarios:} $P(E) = \sum_i P(S_i) \cdot P(E|S_i)$
\end{itemize}

\textbf{ID Conventions:}
\begin{itemize}
    \item Forecast: \texttt{FCT-\{DOMAIN\}-\{YYYY\}-\{NNN\}}
    \item Session: \texttt{EBF-S-\{YYYY\}-\{MM\}-\{DD\}-\{DOMAIN\}-\{NNN\}}
\end{itemize}
\end{tcolorbox}

% =============================================================================
% SECTION 1: INTRODUCTION
% =============================================================================
\section{Introduction}
\label{sec:fct-intro}

\subsection{The Fundamental Question}

\begin{tcolorbox}[colback=yellow!10!white,colframe=yellow!50!black,title=Central Question]
\textbf{How do we systematically track, update, and calibrate probabilistic forecasts within the EBF framework?}
\end{tcolorbox}

% -----------------------------------------------------------------------------
% SCOPE BOX
% -----------------------------------------------------------------------------
\begin{tcolorbox}[colback=orange!5!white,colframe=orange!75!black,title=Appendix Scope]
\textbf{Appendix Scope:} This appendix establishes a systematic methodology for tracking, updating, and evaluating probabilistic forecasts.

\textbf{Ziel:} Systematische Methodik für Forecast-Tracking, -Updating und -Evaluation.

\textbf{In-Scope:}
\begin{itemize}
    \item Forecast Registry architecture and ID conventions
    \item Probability history tracking and Bayesian updating
    \item Calibration metrics (Brier Score, Log Loss)
    \item Conditional scenario analysis
    \item Integration with EBF 10C and $\Psi$-dimensions
\end{itemize}

\textbf{Out-of-Scope:}
\begin{itemize}
    \item Domain-specific models (see MOD-015, MOD-GEO-001)
    \item Parameter estimation (see Appendix AN: METHOD-LLMMC)
    \item General prediction theory (see Appendix S: PREDICT-FALSIFY)
\end{itemize}

\textbf{Deliverables / Lieferobjekte:}
\begin{itemize}
    \item FCT-1 to FCT-5: Five axioms for forecast methodology
    \item Four-database architecture specification
    \item Update and resolution workflows
    \item Worked examples (political referendum, geopolitical event)
\end{itemize}
\end{tcolorbox}

Forecasting is inherently uncertain, but \textit{systematic} forecasting enables learning. The EBF Forecast Protocol addresses three challenges:

\begin{enumerate}
    \item \textbf{Tracking:} Recording predictions with precise resolution criteria
    \item \textbf{Updating:} Bayesian probability revision as new information arrives
    \item \textbf{Calibrating:} Measuring forecast accuracy to improve future models
\end{enumerate}

\subsection{Scope and Domains}

The EBF supports forecasts across six domains:

\begin{table}[h]
\centering
\begin{tabular}{lll}
\toprule
\textbf{Domain} & \textbf{Description} & \textbf{Example} \\
\midrule
POL & Political referenda, elections & SRG Initiative, 10-Million \\
GEO & Geopolitical events, conflicts & Iran attack, Ukraine \\
ECO & Economic forecasts & Recession, inflation \\
FIN & Financial market predictions & Interest rates, FX \\
BEH & Behavioral predictions & Adoption, compliance \\
ORG & Organizational outcomes & Projects, M\&A \\
\bottomrule
\end{tabular}
\caption{EBF Forecast Domains}
\label{tab:forecast-domains}
\end{table}

% =============================================================================
% SECTION 2: AXIOMS
% =============================================================================
\section{Axioms}
\label{sec:fct-axioms}

\begin{axiom}[FCT-1: Resolution Criteria Ex Ante]
\label{ax:fct-1}
Every forecast must define precise resolution criteria \textit{before} the event:
\begin{itemize}
    \item What counts as ``Yes'' (event occurred)?
    \item What counts as ``No'' (event did not occur)?
    \item What actors/actions are included/excluded?
    \item What is the time horizon?
\end{itemize}
\end{axiom}

\begin{axiom}[FCT-2: Probability History]
\label{ax:fct-2}
Every forecast maintains a time-series of probability estimates:
\[
\mathcal{H} = \{(t_i, P_i, \text{CI}_i, R_i)\}_{i=1}^{n}
\]
where $t_i$ is the date, $P_i$ the probability, $\text{CI}_i$ the confidence interval, and $R_i$ the reasoning.
\end{axiom}

\begin{axiom}[FCT-3: Bayesian Updating]
\label{ax:fct-3}
New information triggers probability revision via Bayes' rule:
\[
P(\text{Event}|\text{New Info}) = \frac{P(\text{New Info}|\text{Event}) \cdot P(\text{Event})}{P(\text{New Info})}
\]
Updates must document the prior, likelihood, and posterior explicitly.
\end{axiom}

\begin{axiom}[FCT-4: Calibration Measurement]
\label{ax:fct-4}
After resolution, forecasts are evaluated using:
\begin{itemize}
    \item \textbf{Brier Score:} $BS = (P - O)^2$ where $O \in \{0, 1\}$
    \item \textbf{Log Loss:} $LL = -[O \log(P) + (1-O) \log(1-P)]$
    \item \textbf{Direction Correct:} Did $P > 0.5$ match $O = 1$?
\end{itemize}
\end{axiom}

\begin{axiom}[FCT-5: Four-Database Architecture]
\label{ax:fct-5}
Every forecast must be recorded in four databases:
\begin{enumerate}
    \item \texttt{forecast-registry.yaml} (primary)
    \item \texttt{model-building-session.yaml}
    \item \texttt{model-registry.yaml} (if new model)
    \item \texttt{output-registry.yaml}
\end{enumerate}
Cross-references via \texttt{forecast\_registry\_link} ensure traceability.
\end{axiom}

% =============================================================================
% SECTION 3: METHODOLOGY
% =============================================================================
\section{Methodology}
\label{sec:fct-methodology}

\subsection{Forecast Types}

\subsubsection{Type 1: Fixed-Date Events (Referenda)}

Political referenda have a known resolution date:

\begin{itemize}
    \item \textbf{Question:} Will initiative X pass on date Y?
    \item \textbf{Model:} Referendum Dynamics Model (MOD-015)
    \item \textbf{Data:} Polls, historical corrections, segment analysis
    \item \textbf{Resolution:} Official BFS result
\end{itemize}

\textbf{Example:} SRG-Halbierungsinitiative (FCT-POL-2026-001)
\begin{itemize}
    \item Event date: March 8, 2026
    \item Forecast: 52.1\% NEIN [48.1\%-56.5\%]
    \item P(NEIN) = 62.4\%
\end{itemize}

\subsubsection{Type 2: Time-Horizon Events (Geopolitical)}

Geopolitical events have uncertain timing:

\begin{itemize}
    \item \textbf{Question:} Will event X occur within time horizon T?
    \item \textbf{Model:} Game-theoretic (sequential games, backward induction)
    \item \textbf{Data:} Intelligence, statements, troop movements
    \item \textbf{Resolution:} Event occurrence or horizon expiration
\end{itemize}

\textbf{Example:} Iran Attack (FCT-GEO-2026-001)
\begin{itemize}
    \item Time horizon: End of 2026
    \item Short-term (48h): 18\% [12\%-25\%]
    \item Annual: 35\% [20\%-55\%]
    \item Trigger scenarios with conditional probabilities
\end{itemize}

\subsection{Probability History Structure}

Each forecast maintains a probability history:

\begin{verbatim}
probability_history:
  - date: "2026-01-28"
    probability: 0.18
    confidence_interval: [0.12, 0.25]
    confidence: "MEDIUM"
    reasoning: |
      [Detailed reasoning for this estimate]
    key_drivers:
      positive: ["Driver 1", "Driver 2"]
      negative: ["Driver 3", "Driver 4"]
    trigger_scenarios:
      - trigger: "Event A"
        p_given_trigger: 0.60
    updated_by: "Session ID or Update Reason"
\end{verbatim}

\subsection{Conditional Scenarios}

For complex forecasts, decompose into conditional probabilities:

\[
P(\text{Attack}) = \sum_i P(\text{Scenario}_i) \cdot P(\text{Attack}|\text{Scenario}_i)
\]

\textbf{Example:}
\begin{table}[h]
\centering
\begin{tabular}{llll}
\toprule
\textbf{Scenario} & $P(\text{Scenario})$ & $P(\text{Attack}|\text{Scenario})$ & \textbf{Combined} \\
\midrule
Iran Breakout & 0.40 & 0.70 & 0.28 \\
Status Quo & 0.50 & 0.15 & 0.075 \\
Diplomatic Deal & 0.10 & 0.02 & 0.002 \\
\midrule
\textbf{Total} & 1.00 & - & \textbf{0.357} \\
\bottomrule
\end{tabular}
\caption{Conditional Scenario Analysis}
\label{tab:conditional-scenarios}
\end{table}

% =============================================================================
% SECTION 4: CALIBRATION
% =============================================================================
\section{Calibration}
\label{sec:fct-calibration}

\subsection{Brier Score}

The Brier Score measures forecast accuracy:

\[
BS = \frac{1}{N} \sum_{i=1}^{N} (P_i - O_i)^2
\]

\begin{itemize}
    \item $BS = 0$: Perfect forecast
    \item $BS = 0.25$: Random guessing (for 50/50 events)
    \item $BS = 1$: Worst possible forecast
\end{itemize}

\subsection{Log Loss}

Log loss penalizes overconfident wrong forecasts:

\[
LL = -\frac{1}{N} \sum_{i=1}^{N} [O_i \log(P_i) + (1-O_i) \log(1-P_i)]
\]

\subsection{Calibration Curve}

A well-calibrated forecaster has:
\[
\text{Fraction Correct}(P) \approx P
\]

For example, events forecasted at 70\% should occur about 70\% of the time.

% =============================================================================
% SECTION 5: DATABASE ARCHITECTURE
% =============================================================================
\section{Database Architecture}
\label{sec:fct-database}

\subsection{Four-Database Model}

\begin{figure}[h]
\centering
\begin{verbatim}
                    ┌─────────────────────────┐
                    │  1. FORECAST-REGISTRY   │
                    │  FCT-XXX-YYYY-NNN       │
                    │  (Primary)              │
                    └───────────┬─────────────┘
                                │
         ┌──────────────────────┼──────────────────────┐
         │                      │                      │
         ▼                      ▼                      ▼
  ┌─────────────┐      ┌─────────────┐      ┌─────────────┐
  │ 2. SESSION  │      │ 3. MODEL    │      │ 4. OUTPUT   │
  │ EBF-S-...   │      │ MOD-XXX-NNN │      │ EBF-OUT-... │
  └─────────────┘      └─────────────┘      └─────────────┘
\end{verbatim}
\caption{Four-Database Architecture}
\label{fig:four-db}
\end{figure}

\subsection{Cross-References}

All databases link via \texttt{forecast\_registry\_link}:

\begin{verbatim}
# In model-building-session.yaml:
forecast_registry_link: "FCT-GEO-2026-001"

# In output-registry.yaml:
forecast_registry_link: "FCT-GEO-2026-001"
\end{verbatim}

\subsection{ID Conventions}

\begin{table}[h]
\centering
\begin{tabular}{lll}
\toprule
\textbf{Type} & \textbf{Format} & \textbf{Example} \\
\midrule
Forecast & FCT-\{DOMAIN\}-\{YYYY\}-\{NNN\} & FCT-GEO-2026-001 \\
Session & EBF-S-\{YYYY\}-\{MM\}-\{DD\}-\{DOMAIN\}-\{NNN\} & EBF-S-2026-01-28-GEO-001 \\
Model & MOD-\{DOMAIN\}-\{NNN\} & MOD-GEO-001 \\
Output & EBF-OUT-\{DOMAIN\}-\{NNN\} & EBF-OUT-GEO-001 \\
\bottomrule
\end{tabular}
\caption{ID Conventions}
\label{tab:id-conventions}
\end{table}

% =============================================================================
% SECTION 6: WORKED EXAMPLES
% =============================================================================
\section{Worked Examples}
\label{sec:fct-examples}

\subsection{Example 1: Political Referendum (SRG-Halbierung)}

\textbf{Setup:}
\begin{itemize}
    \item Question: Will the SRG initiative pass on March 8, 2026?
    \item Model: Referendum Dynamics Model Extended (MOD-015)
    \item Data: gfs.bern, Tamedia, Sotomo polls
\end{itemize}

\textbf{Analysis:}
\begin{enumerate}
    \item Weighted baseline from polls: 49.2\% NEIN, 48.4\% JA
    \item Historical correction: -7pp for JA (No-Billag, RTVG, Medienpaket)
    \item Mobilization asymmetry: +3pp for NEIN
    \item Monte Carlo simulation: N=10,000
\end{enumerate}

\textbf{Result:}
\[
P(\text{NEIN}) = 52.1\% \quad [48.1\%, 56.5\%]
\]
\[
P(\text{Ablehnung}) = 62.4\%
\]

\subsection{Example 2: Geopolitical Event (Iran Attack)}

\textbf{Setup:}
\begin{itemize}
    \item Question: Will US/Israel attack Iranian nuclear facilities in 2026?
    \item Model: Sequential Game with Backward Induction (MOD-GEO-001)
    \item Data: Troop movements, statements, IAEA reports
\end{itemize}

\textbf{Analysis:}
\begin{enumerate}
    \item Game-theoretic equilibrium: Negotiate $\rightarrow$ Deal
    \item Short-term (48h): P = 18\%
    \item Trigger scenarios: Proxy attack (+42pp), Breakout (+52pp)
    \item Annual cumulative: P = 35\%
\end{enumerate}

\textbf{Result:}
\[
P(\text{Attack}|48h) = 18\% \quad [12\%, 25\%]
\]
\[
P(\text{Attack}|2026) = 35\% \quad [20\%, 55\%]
\]

% =============================================================================
% SECTION 7: UPDATE WORKFLOW
% =============================================================================
\section{Update Workflow}
\label{sec:fct-update}

\subsection{When to Update}

\begin{itemize}
    \item New polls (political forecasts)
    \item Significant events (geopolitical)
    \item IAEA reports, troop movements
    \item High-level statements
    \item Time decay (approaching deadline)
\end{itemize}

\subsection{Update Protocol}

\begin{enumerate}
    \item Add new entry to \texttt{probability\_history}
    \item Document reasoning (what changed?)
    \item Update \texttt{key\_drivers}
    \item Revise \texttt{conditional\_scenarios} if needed
    \item Update \texttt{forecast\_updated} date
    \item Commit with message: \texttt{update(forecast): FCT-XXX reason}
\end{enumerate}

% =============================================================================
% SECTION 8: RESOLUTION WORKFLOW
% =============================================================================
\section{Resolution Workflow}
\label{sec:fct-resolution}

\subsection{After Event Occurs}

\begin{verbatim}
actual_outcome:
  resolved: true
  date: "2026-03-08"
  result: "nein"  # or "ja", "attack", "no_attack"
  details: "52.1% NEIN, 44.9% JA"
  source: "BFS"

calibration:
  brier_score: 0.001  # (0.521 - 0.521)²
  log_loss: 0.65
  direction_correct: true
  within_ci: true
\end{verbatim}

\subsection{Aggregate Calibration}

Update \texttt{calibration\_summary} in forecast-registry.yaml:

\begin{verbatim}
calibration_summary:
  total_forecasts: 3
  resolved_forecasts: 1
  pending_forecasts: 2
  metrics:
    brier_score: 0.001
    avg_log_loss: 0.65
  by_domain:
    POL:
      count: 2
      resolved: 1
      avg_brier: 0.001
\end{verbatim}

% =============================================================================
% SECTION 9: INTEGRATION WITH EBF
% =============================================================================
\section{Integration with EBF Framework}
\label{sec:fct-ebf}

\subsection{10C Mapping}

Forecasts map to 10C dimensions:

\begin{itemize}
    \item \textbf{WHO:} Actors, segments, stakeholders
    \item \textbf{WHAT:} Utility dimensions at stake
    \item \textbf{HOW:} Interaction structure (game theory)
    \item \textbf{WHEN:} Context dimensions ($\Psi$)
    \item \textbf{WHERE:} Parameter sources (BBB)
    \item \textbf{AWARE:} Information, salience
    \item \textbf{READY:} Willingness thresholds
    \item \textbf{STAGE:} Decision stage in journey
\end{itemize}

\subsection{$\Psi$-Dimensions in Forecasting}

Context dimensions affect forecast parameters:

\begin{itemize}
    \item $\Psi_I$ (Institutional): Voting rules, defaults
    \item $\Psi_S$ (Social): Social pressure, norms
    \item $\Psi_K$ (Cultural): Values, traditions
    \item $\Psi_T$ (Temporal): Time pressure, deadlines
    \item $\Psi_E$ (Economic): Resources, constraints
\end{itemize}

% =============================================================================
% SECTION 10: RESULTS
% =============================================================================
\section{Results}
\label{sec:fct-results}

\subsection{Current Forecast Portfolio}

As of January 2026, the EBF Forecast Registry contains:

\begin{table}[h]
\centering
\begin{tabular}{llllr}
\toprule
\textbf{ID} & \textbf{Domain} & \textbf{Event} & \textbf{Status} & \textbf{P(Event)} \\
\midrule
FCT-POL-2025-001 & POL & 10-Million Initiative & PENDING & 48\% JA \\
FCT-POL-2026-001 & POL & SRG-Halbierung & PENDING & 37.6\% JA \\
FCT-GEO-2026-001 & GEO & Iran Attack (48h) & PENDING & 18\% \\
\bottomrule
\end{tabular}
\caption{Current Forecast Portfolio (January 2026)}
\label{tab:forecast-portfolio}
\end{table}

\subsection{Calibration Results}

No forecasts have been resolved yet. After resolution, calibration metrics will be computed:

\begin{itemize}
    \item \textbf{Aggregate Brier Score:} Target $BS < 0.15$
    \item \textbf{Calibration Curve:} Target $|\text{Predicted} - \text{Actual}| < 5\%$ per decile
    \item \textbf{Direction Accuracy:} Target $> 80\%$ correct direction calls
\end{itemize}

\subsection{Key Findings}

\begin{enumerate}
    \item \textbf{Political Referenda:} Historical corrections are significant ($-7$pp for media initiatives)
    \item \textbf{Geopolitical Events:} Trigger scenarios drive most of the variance
    \item \textbf{Four-Database Architecture:} Cross-referencing ensures traceability
\end{enumerate}

% =============================================================================
% SECTION 11: IMPLICATIONS
% =============================================================================
\section{Implications}
\label{sec:fct-implications}

\subsection{For EBF Practice}

\begin{enumerate}
    \item \textbf{Every model produces forecasts:} The model-building workflow (Appendix EEE) should always generate testable predictions
    \item \textbf{Forecasts must be tracked:} Ad-hoc predictions without registry entries violate FCT-5
    \item \textbf{Updates are learning opportunities:} Each probability update documents what changed and why
    \item \textbf{Calibration drives improvement:} Post-resolution analysis identifies systematic biases
\end{enumerate}

\subsection{For Decision-Making}

\begin{enumerate}
    \item \textbf{Probability ranges matter:} Point estimates are less useful than confidence intervals
    \item \textbf{Conditional scenarios enable planning:} Decomposition into $P(S_i) \cdot P(E|S_i)$ supports contingency planning
    \item \textbf{Trigger events deserve monitoring:} High-impact triggers (e.g., proxy attack: $+42$pp) should be tracked
\end{enumerate}

\subsection{For Research}

\begin{enumerate}
    \item \textbf{Transparent methodology:} All assumptions, models, and data sources are documented
    \item \textbf{Reproducible analysis:} Probability histories enable replication
    \item \textbf{Cumulative learning:} Calibration metrics aggregate across forecasts
\end{enumerate}

% =============================================================================
% SECTION 12: SUMMARY
% =============================================================================
\section{Summary}
\label{sec:fct-summary}

\begin{tcolorbox}[colback=blue!5!white,colframe=blue!75!black,title=Summary]
\textbf{Core Contributions:}
\begin{enumerate}
    \item \textbf{Five Axioms (FCT-1 to FCT-5):} Formal requirements for EBF-compliant forecasts
    \item \textbf{Four-Database Architecture:} Systematic tracking across forecast, session, model, and output registries
    \item \textbf{Probability History:} Time-series tracking with Bayesian updating
    \item \textbf{Calibration Framework:} Brier Score, Log Loss, and calibration curves
    \item \textbf{Workflows:} Create, update, and resolve protocols
\end{enumerate}

\textbf{Key Takeaways:}
\begin{itemize}
    \item Forecasts are \textit{systematically tracked}, not ad-hoc predictions
    \item Resolution criteria are defined \textit{ex ante}, not post-hoc
    \item Calibration enables \textit{learning and improvement} over time
    \item Integration with 10C and $\Psi$-dimensions grounds forecasts in behavioral economics
\end{itemize}
\end{tcolorbox}

% =============================================================================
% SECTION 13: CRITICAL FOUNDATIONS
% =============================================================================
\section{Critical Foundations}
\label{sec:fct-foundations}

This appendix builds on foundational work in:

\begin{enumerate}
    \item \textbf{Proper Scoring Rules} \citep{brier1950verification, good1952rational}: Mathematical foundations for measuring forecast accuracy
    \item \textbf{Superforecasting} \citep{tetlock2015superforecasting}: Empirical evidence that calibrated forecasters outperform experts
    \item \textbf{Bayesian Updating} \citep{bayes1763essay}: Probability revision with new information
    \item \textbf{Game Theory} \citep{selten1965spieltheoretische}: Strategic interaction modeling for geopolitical forecasts
    \item \textbf{EBF Context Framework} (Appendix V): $\Psi$-dimensions affect forecast parameters
\end{enumerate}

\subsection{Theoretical Grounding}

The forecast methodology is grounded in three principles:

\begin{description}
    \item[Coherence:] Probabilities must satisfy the Kolmogorov axioms
    \item[Calibration:] Predicted probabilities should match actual frequencies
    \item[Resolution:] Forecasts should discriminate between events and non-events
\end{description}

% =============================================================================
% SECTION 14: OPEN ISSUES
% =============================================================================
\section{Open Issues}
\label{sec:fct-open-issues}

\begin{enumerate}
    \item \textbf{Small Sample Size:} With only 3 forecasts, calibration metrics are not yet statistically meaningful

    \item \textbf{Domain-Specific Calibration:} Different domains (POL, GEO, ECO) may require different calibration standards

    \item \textbf{Trigger Scenario Validation:} Conditional probabilities $P(E|S_i)$ are difficult to validate empirically

    \item \textbf{Model Selection:} Which model to use for which domain? (see Appendix EEE for guidance)

    \item \textbf{Automation:} Should probability updates be triggered automatically by news events?

    \item \textbf{Ensemble Forecasting:} Should we combine multiple models for the same forecast?

    \item \textbf{Long-Horizon Forecasts:} How to handle forecasts with $T > 1$ year?
\end{enumerate}

\subsection{Future Work}

\begin{itemize}
    \item Implement automated calibration dashboards
    \item Develop domain-specific scoring rules
    \item Integrate with external prediction markets for validation
    \item Build ensemble models combining multiple forecasting approaches
\end{itemize}

% =============================================================================
% REFERENCES
% =============================================================================
\section{References}
\label{sec:fct-references}

\nocite{bcm_master}

Key references for forecast methodology:
\begin{itemize}
    \item Tetlock, P. E. (2005). Expert Political Judgment
    \item Tetlock, P. E. \& Gardner, D. (2015). Superforecasting
    \item Brier, G. W. (1950). Verification of Forecasts
    \item Good, I. J. (1952). Rational Decisions (Log Loss)
    \item Gneiting, T. \& Raftery, A. E. (2007). Strictly Proper Scoring Rules
\end{itemize}

% =============================================================================
% GLOSSARY
% =============================================================================
\section{Glossary}
\label{sec:fct-glossary}

See Appendix G (Glossary) for comprehensive definitions. Key terms:

\begin{description}
    \item[Brier Score] Quadratic scoring rule: $(P - O)^2$
    \item[Calibration] Alignment between predicted probabilities and actual frequencies
    \item[Log Loss] Logarithmic scoring rule penalizing overconfidence
    \item[Probability History] Time-series of probability estimates for a forecast
    \item[Resolution Criteria] Ex-ante definition of what counts as event occurrence
\end{description}

% =============================================================================
% CHANGELOG
% =============================================================================
\section*{Changelog}

\begin{tabular}{lll}
\textbf{Version} & \textbf{Date} & \textbf{Changes} \\
\hline
1.0 & 2026-01-30 & Initial version \\
\end{tabular}

\end{document}
